\documentclass[11pt]{article}

\usepackage[margin=1in]{geometry}
\usepackage{amsmath,amssymb}
\usepackage{booktabs}
\usepackage{array}
\usepackage{hyperref}
\usepackage{xcolor}

\hypersetup{
  colorlinks=true,
  linkcolor=blue,
  urlcolor=blue
}

\title{\textbf{Counterfactual LogFC Conversion for Residual-Scale DEDML Estimates}}
\author{Prepared for Yuxin Yin}
\date{April 30, 2026}

\begin{document}
\maketitle

\section*{Motivation}

DEDML estimates treatment effects after outcome residualization and treatment
residualization. Its primary coefficient,
\[
  \widehat{\beta}_{gc}^{\mathrm{res}},
\]
is therefore on a standardized residual scale for gene \(g\) and cell type
\(c\), not directly on the natural-log fold-change scale.

The earlier post-hoc conversion used the fitted nuisance mean
\(\widehat m(X)=\widehat E(Y\mid X)\) as if it were the untreated baseline
mean. This is not generally correct. If treatment is associated with
confounders, then
\[
  m(X)
  =
  E(Y\mid X)
  =
  \mu_0(X)\{1+e(X)(\exp(\tau)-1)\},
\]
where
\[
  e(X)=P(D=1\mid X),
  \qquad
  \mu_0(X)=E(Y\mid D=0,X).
\]
Thus \(m(X)\) is a treatment-prevalence-weighted mixture, not the control
counterfactual mean. The reported logFC should target \(\tau\), the effect
relative to \(\mu_0(X)\).

\section*{New Conversion}

Let
\[
  \delta=\exp(\tau)-1.
\]
Because
\[
  m_i = \mu_{0i}(1+\widehat e_i\delta),
\]
we can write the untreated baseline as
\[
  \mu_{0i} = \frac{m_i}{1+\widehat e_i\delta}.
\]

For Poisson-scale residuals, the residualized mean shift induced by \(\delta\)
is approximately
\[
  \frac{\mu_{0i}\delta}{\sqrt{m_i}}
  =
  \frac{\sqrt{m_i}\delta}{1+\widehat e_i\delta}.
\]
Therefore the corrected conversion solves the one-dimensional equation
\[
  \widehat{\beta}_{gc}^{\mathrm{res}}
  =
  \frac{1}{n_{gc}}
  \sum_{i:C_i=c}
  \frac{\sqrt{\widehat m_{ig}}\delta}
       {1+\widehat e_i\delta},
\]
and then reports
\[
  \widehat\tau_{gc} = \log(1+\widehat\delta_{gc}).
\]

For negative-binomial residuals with size parameter \(\widehat\phi_g\), the
same idea is used with the NB residual scale:
\[
  \widehat{\beta}_{gc}^{\mathrm{res}}
  =
  \frac{1}{n_{gc}}
  \sum_{i:C_i=c}
  \frac{
    \widehat m_{ig}\delta
  }{
    \{1+\widehat e_i\delta\}
    \sqrt{\widehat m_{ig}+\widehat m_{ig}^2/\widehat\phi_g}
  }.
\]
The solution again gives
\[
  \widehat\tau_{gc} = \log(1+\widehat\delta_{gc}).
\]

The inference remains unchanged. DEDML p-values and test statistics continue
to use the residual-scale regression. The new conversion only changes the
reported effect-size column \texttt{estimate\_logfc}.

\section*{Implementation}

The DEDML package now stores the cross-fitted treatment propensity
\(\widehat e_i\) alongside the outcome residual quantities. For each
gene-celltype result row, \texttt{estimate\_logfc} is computed by numerically
solving the equation above. The output also includes summary diagnostics:
\[
  \texttt{muhat\_mean},
  \qquad
  \texttt{muhat\_factor},
  \qquad
  \texttt{muhat\_propensity\_mean},
  \qquad
  \texttt{muhat\_nb\_size}.
\]

The package version was bumped to \texttt{0.1.2}.

\section*{Benchmark Design}

We benchmarked the conversion on the observed-confounding synthetic
single-cell setting:
\begin{itemize}
  \item donors: 50;
  \item genes: 500;
  \item cell types: 6;
  \item repeated samples per donor: none;
  \item mean cells per sample: 700;
  \item DE gene fraction: 8\%;
  \item DE cell types per DE gene: 2;
  \item simulated natural-log fold-change magnitude: around 0.45;
  \item DEDML folds: 3;
  \item cores: 20;
  \item donor model: logistic GLM;
  \item outcome model: Poisson GLM with batched exact IRLS.
\end{itemize}

The treatment model produced nontrivial estimated propensities:
\[
\begin{array}{lr}
\text{donors} & 50 \\
\text{treated donors} & 24 \\
\min(\widehat e) & 0.0001 \\
\max(\widehat e) & 0.9999 \\
\operatorname{mean}(\widehat e) & 0.458 \\
\operatorname{sd}(\widehat e) & 0.499 \\
\operatorname{mean}(\widehat e\mid D=1) & 0.577 \\
\operatorname{mean}(\widehat e\mid D=0) & 0.349
\end{array}
\]

\section*{Algebraic Check}

First, we tested the conversion directly under known residual-scale effects.
The new conversion recovers the true untreated-baseline logFC to numerical
precision, while the old mixed-mean conversion can be severely biased.

\begin{center}
\begin{tabular}{lrrrr}
\toprule
Distribution & Old mean abs. error & Old max abs. error & New mean abs. error & New max abs. error \\
\midrule
Poisson & 0.423 & 2.938 & \(8.3\times10^{-10}\) & \(3.2\times10^{-9}\) \\
NB      & 0.445 & 3.196 & \(8.2\times10^{-10}\) & \(3.3\times10^{-9}\) \\
\bottomrule
\end{tabular}
\end{center}

\section*{End-to-End DEDML Results}

The table below compares the old mixed-mean conversion with the new
untreated-baseline conversion against the simulated truth. The hypothesis test
is unchanged, so FDR and power are not affected by this conversion.

\begin{center}
\begin{tabular}{lrrrrrr}
\toprule
Conversion & Kendall DE & Spearman DE & Pearson DE & Bias DE & MAE DE & RMSE DE \\
\midrule
Old mixed mean & 0.572 & 0.794 & 0.891 & -0.048 & 0.215 & 0.250 \\
New untreated baseline & 0.571 & 0.793 & 0.896 & -0.006 & 0.225 & 0.253 \\
\bottomrule
\end{tabular}
\end{center}

The main empirical improvement is bias reduction: the old conversion
underestimates DE logFC by about \(-0.048\), whereas the new conversion reduces
this bias to about \(-0.006\). Correlations are similar because, in the
end-to-end benchmark, much of the remaining error comes from the residual-scale
DEDML estimate itself rather than from the scale conversion.

\section*{Interpretation}

The new conversion fixes the target mismatch. It reports an effect size on the
untreated-baseline logFC scale using only estimable quantities:
\[
  \widehat m(X),
  \qquad
  \widehat e(X),
  \qquad
  \widehat\phi_g \text{ for NB}.
\]
It does not use simulation truth and therefore is applicable to real data.

The remaining caveat is that \texttt{estimate\_logfc} is still a post-hoc
effect-size interpretation of a residual-scale DEDML coefficient. It should be
reported together with the residual-scale estimate, statistic, and p-value.

\end{document}
